% =====================================================================
%  CSE 344 - System Programming - Homework #6 : "Magic Academy"
%  Report template (Overleaf ready).
%
%  RECOMMENDED COMPILER: LuaLaTeX or XeLaTeX  (needed for the "Caveat"
%  hand-written font via fontspec).  It also compiles under pdfLaTeX,
%  in which case the hand-written annotations fall back to an italic font.
%
%  In Overleaf:  Menu  ->  Compiler  ->  LuaLaTeX
% =====================================================================
% !TeX program = lualatex
\documentclass[11pt,a4paper]{article}

\usepackage[T1]{fontenc}
\usepackage[utf8]{inputenc}        % harmless under LuaLaTeX/XeLaTeX
\usepackage[margin=2.4cm]{geometry}
\usepackage{microtype}
\usepackage{graphicx}
\usepackage{xcolor}
\usepackage{hyperref}
\usepackage{enumitem}
\usepackage{listings}
\usepackage{booktabs}
\usepackage{tcolorbox}
\usepackage{tikz}
\usetikzlibrary{positioning,arrows.meta,shapes.multipart,fit,backgrounds,calc}

% ---------------------------------------------------------------------
%  Hand-written "Caveat" font for stylized annotations / blockquotes.
%  Uses fontspec when available, otherwise degrades gracefully.
% ---------------------------------------------------------------------
\usepackage{iftex}
\ifPDFTeX
  % pdfLaTeX fallback: an italic serif stands in for the hand-written look.
  \newcommand{\handfont}{\itshape}
\else
  \usepackage{fontspec}
  % "Caveat" ships with most Overleaf installations; if missing, swap the
  % name below for any installed hand-written font (e.g. "Patrick Hand").
  \newfontfamily{\CaveatFamily}{Caveat}
  \newcommand{\handfont}{\CaveatFamily}
\fi

% A margin/annotation note that looks hand-written.
\newcommand{\annot}[1]{{\handfont\large\color{teal}#1}}

% A "scribbled remark" callout box.
\newtcolorbox{remark}{
  colback=teal!4, colframe=teal!55!black, boxrule=0.6pt, arc=2pt,
  left=8pt, right=8pt, top=6pt, bottom=6pt,
  fontupper=\handfont\large
}

% ---------------------------------------------------------------------
%  Source-code listing style
% ---------------------------------------------------------------------
\lstdefinestyle{c}{
  language=C, basicstyle=\ttfamily\footnotesize, breaklines=true,
  keywordstyle=\color{blue!70!black}, commentstyle=\color{green!45!black},
  stringstyle=\color{purple!70!black}, numbers=left, numberstyle=\tiny\color{gray},
  frame=single, framesep=4pt, rulecolor=\color{gray!40},
  showstringspaces=false, tabsize=4, xleftmargin=14pt
}
\lstset{style=c}

\hypersetup{colorlinks=true, linkcolor=blue!50!black, urlcolor=blue!55!black}

% =====================================================================
\begin{document}

% ----------------------------- Title page ----------------------------
\begin{titlepage}
\centering
\vspace*{2.5cm}
{\Huge\bfseries Magic Academy\par}
\vspace{0.4cm}
{\Large A Single-Threaded \texttt{select()} TCP Server\par}
\vspace{1.6cm}
{\large CSE 344 -- System Programming\par}
{\large Homework \#6\par}
\vspace{2.2cm}
\begin{tabular}{rl}
  \textbf{Name:}        & \texttt{<Your Name>}\\[4pt]
  \textbf{Student ID:}  & \texttt{<Your ID>}\\[4pt]
  \textbf{Course:}      & CSE 344 -- System Programming\\[4pt]
  \textbf{Homework:}    & \#6 -- Magic Academy\\[4pt]
  \textbf{Date:}        & \today\\
\end{tabular}
\vfill
\annot{Hogwarts opens its gates every morning\ldots\ and closes them on \texttt{Ctrl+C}.}
\vfill
\end{titlepage}

\tableofcontents
\newpage

% =====================================================================
\section{Overview}
This homework implements three programs that communicate over TCP:

\begin{itemize}[nosep]
  \item \textbf{\texttt{hogwarts}} -- the academy server. It loads a list of
        magical ingredients, accepts wizard and professor connections, keeps
        per-client state (username, type, line buffer, spellbook, last-activity
        timestamp), enforces a maximum number of simultaneous clients and an
        idle timeout, and shuts down cleanly on \texttt{SIGINT}.
  \item \textbf{\texttt{wizard}} -- a TCP client that auto-enrolls as a wizard
        and can issue \texttt{BREW}, \texttt{CONSUME}, \texttt{SPELLBOOK},
        \texttt{INSPECT}, \texttt{SCROLL}, \texttt{ROSTER}, \texttt{APPARATE}.
  \item \textbf{\texttt{professor}} -- a TCP client that auto-enrolls as a
        professor; it may query (\texttt{INSPECT}, \texttt{SCROLL},
        \texttt{ROSTER}) but the server rejects \texttt{BREW}/\texttt{CONSUME}
        from it with \texttt{ERR UNAUTHORIZED}.
\end{itemize}

\begin{remark}
The whole server is one process, one thread, one \texttt{select()} loop --
no \texttt{pthread}, no \texttt{fork} per client, no \texttt{sleep()}.
\end{remark}

\paragraph{Build \& run.}
\begin{lstlisting}[language=bash]
$ make                       # builds hogwarts, wizard, professor
$ ./hogwarts -p 5000 -s ingredients.txt -l server.log -n 4 -t 30
$ ./wizard    127.0.0.1 5000 hermione
$ ./professor 127.0.0.1 5000 dumbledore
$ make clean
\end{lstlisting}

% =====================================================================
\section{System Design}

\subsection{One loop, many clients}
The server keeps a fixed-size array \texttt{client\_t clients[max\_clients]}.
A free slot is marked by \texttt{fd == -1}. Each iteration of the main loop:

\begin{enumerate}[nosep]
  \item \textbf{Idle sweep.} Walk the array; for every connected client compute
        \texttt{elapsed = now - last\_active}. If \texttt{elapsed >= TIMEOUT},
        send \texttt{TIMEOUT DISCONNECT}, log \texttt{TIMEOUT} and
        \texttt{CLIENT\_DISCONNECTED reason=timeout}, \texttt{close()} the fd
        and free the slot.
  \item \textbf{Build the \texttt{fd\_set}.} \texttt{FD\_ZERO}, then
        \texttt{FD\_SET} the listening socket and every connected client fd;
        track \texttt{maxfd}.
  \item \textbf{Recompute the timeout.} For every connected client compute the
        remaining time \texttt{TIMEOUT - elapsed}; the smallest value becomes
        the \texttt{struct timeval} passed to \texttt{select()}. If there are
        no clients, \texttt{NULL} is passed (block indefinitely).
  \item \textbf{\texttt{select()}.} Wait for readability on any fd, or until the
        nearest idle deadline expires.
  \item \textbf{Dispatch.} If the listening socket is ready, \texttt{accept()} a
        new connection (\S\ref{sec:capacity}). For every ready client fd, read
        available bytes into its line buffer and process complete lines.
\end{enumerate}

\annot{The select() timeout is recomputed from scratch \emph{before every call}
-- that is how idle disconnection happens without a timer thread or sleep().}

\subsection{Idle timeout without threads or sleep}
\texttt{time()} (whole seconds, monotone enough for a 1\,s-granularity timeout)
stamps \texttt{last\_active} on every byte received from a client. Because the
sweep runs \emph{before} \texttt{select()} and the \texttt{select()} timeout is
the minimum remaining lifetime, the loop always wakes up no later than the next
client's deadline, disconnects it, and recomputes. \texttt{SIGINT} is handled by
not setting \texttt{SA\_RESTART}, so \texttt{select()} returns \texttt{EINTR} and
the loop notices the shutdown flag promptly.

\subsection{Capacity enforcement}\label{sec:capacity}
On a new connection the server always calls \texttt{accept()} first (so the
client gets a real error rather than a silent hang). If the number of occupied
slots already equals \texttt{max\_clients}, it sends \texttt{ERR HOGWARTS\_FULL},
logs \texttt{REJECTED fd=N ip=X.X.X.X reason=HOGWARTS\_FULL}, and
\texttt{close()}s the fd \emph{without} storing it. A slot is reclaimed the
instant a client disconnects (APPARATE / hangup / timeout), so capacity
recovers immediately.

\subsection{Block diagram}
\begin{center}
\begin{tikzpicture}[
  font=\small,
  box/.style={draw, rounded corners=3pt, align=center, minimum height=10mm,
              inner sep=5pt, fill=blue!4},
  proc/.style={draw, rounded corners=8pt, align=center, fill=teal!8,
               minimum height=12mm, inner sep=6pt},
  store/.style={draw, align=center, fill=orange!8, minimum height=10mm,
                inner sep=5pt},
  rec/.style={draw, rectangle split, rectangle split parts=2, align=left,
              fill=purple!5, inner sep=4pt},
  flow/.style={-{Stealth[length=2.4mm]}, thick},
  thin flow/.style={-{Stealth[length=2mm]}, gray}
]

% --- the kernel side ---
\node[box] (listen) {listening socket\\ \texttt{listen\_fd}};
\node[box, below=10mm of listen] (c1) {client socket fd$_1$};
\node[box, below=6mm of c1]      (c2) {client socket fd$_2$};
\node[box, below=6mm of c2]      (cn) {client socket fd$_n$};

% --- the select() core ---
\node[proc, right=26mm of c1, yshift=4mm] (sel)
  {\texttt{select(maxfd+1,\&rfds,\dots,\&tv)}\\[2pt]
   \scriptsize tv = min remaining idle time};
\node[proc, above=9mm of sel] (loop)
  {single-threaded main loop\\ \scriptsize sweep $\to$ build fd\_set $\to$ select $\to$ dispatch};

% --- per-client state ---
\node[rec, right=26mm of sel] (state)
  {\texttt{client\_t} \nodepart{two}
   \scriptsize fd\\ username[32]\\ type[16]\\ line\_buf[512]\\ last\_active\\ spellbook[\,]};
\node[store, above=9mm of state] (clients) {\texttt{clients[max\_clients]}\\ \scriptsize array of slots};
\node[store, below=10mm of state] (ing)   {\texttt{ingredients[\,]}\\ \scriptsize name $\to$ global qty};
\node[store, below=8mm of ing]   (log)    {logfile + stdout\\ \scriptsize timestamped events};

% --- wires ---
\draw[flow] (listen.east) -- node[above,scriptsize]{accept()} (loop.west|-listen) -- (loop.west);
\draw[flow] (c1.east) -- (sel.west);
\draw[flow] (c2.east) -- (sel.west);
\draw[flow] (cn.east) -- (sel.west);
\draw[flow] (loop) -- (sel);
\draw[flow] (sel.east) -- node[above,scriptsize]{read / write line} (state.west);
\draw[thin flow] (clients) -- (state);
\draw[thin flow] (state) -- (ing) node[midway,right,scriptsize]{BREW / CONSUME};
\draw[thin flow] (loop.east) to[out=0,in=90] (clients.north);
\draw[thin flow] (state) -- (log);

\begin{scope}[on background layer]
  \node[draw=gray!40, dashed, rounded corners, fit=(listen)(c1)(c2)(cn),
        label=above:{\scriptsize kernel TCP sockets}] {};
  \node[draw=gray!40, dashed, rounded corners, fit=(loop)(sel),
        label=above:{\scriptsize \texttt{hogwarts} -- one thread}] {};
  \node[draw=gray!40, dashed, rounded corners, fit=(clients)(state)(ing)(log),
        label=above:{\scriptsize server state}] {};
\end{scope}
\end{tikzpicture}
\end{center}

% =====================================================================
\section{Data Structures}

\subsection{\texttt{client\_t}}
\begin{lstlisting}
typedef struct {
    int    fd;                          /* -1 == free slot                     */
    char   username[MAX_USERNAME];      /* from ENROLL                          */
    char   type[MAX_TYPE];              /* "WIZARD" / "PROFESSOR" / ""          */
    int    enrolled;                    /* 1 after a valid ENROLL               */
    char   line_buf[MAX_LINE];          /* partial-read accumulator (<= 512 B)  */
    size_t buf_len;                     /* bytes currently held in line_buf     */
    int    overlong;                    /* 1 while draining an over-long line   */
    time_t last_active;                 /* timestamp of last byte received      */
    long   spellbook[MAX_INGREDIENTS];  /* WIZARD only: per-ingredient holdings */
} client_t;
\end{lstlisting}
All slots live in one heap array \texttt{client\_t *clients = calloc(max\_clients,
sizeof(client\_t))}. Lookup, capacity counting and username-uniqueness checks
are simple linear scans -- \texttt{max\_clients} is small, so this is fine and
keeps the code branch-free of dynamic containers.

\subsection{Spellbook}
A wizard's spellbook is a plain \texttt{long[MAX\_INGREDIENTS]} indexed by the
ingredient's position in the global table. \texttt{BREW} adds to both
\texttt{ingredients[i].qty} and \texttt{spellbook[i]}; \texttt{CONSUME} subtracts
from both. \texttt{CONSUME} is refused with \texttt{ERR INSUFFICIENT\_INGREDIENTS}
unless the wizard holds at least the requested amount (and the global stock does
too). Because every brew first increases the global stock, the global quantity
is always $\ge$ any single wizard's holding, so the global stock can never go
negative. The spellbook is \emph{not} persisted: when the client disconnects the
slot (and the array) is zeroed.

\subsection{Ingredient table}
\begin{lstlisting}
typedef struct { char name[MAX_ING_NAME]; long qty; } ingredient_t;
static ingredient_t g_ing[MAX_INGREDIENTS];
static int          g_ning;
\end{lstlisting}
Loaded once at start-up from \texttt{ingredients.txt}. Each line is parsed with
\texttt{sscanf("\%63s \%ld", \dots)}; lines whose name is not 1--16 uppercase
alphanumerics/underscores, whose quantity is not a positive integer, or which
duplicate an earlier name, are skipped silently.

% =====================================================================
\section{Partial Read / Write Handling}
TCP is a byte stream: one \texttt{read()} may return half a command, one and a
half commands, or several commands glued together.

\paragraph{Reads.} Each client owns \texttt{line\_buf[512]}. On a readable fd the
server does one \texttt{read()} into a scratch buffer and then walks the bytes:
\begin{itemize}[nosep]
  \item a \verb|'\n'| ends a line -- the bytes accumulated in \texttt{line\_buf}
        (with a trailing \verb|'\r'| trimmed) become a NUL-terminated command
        that is dispatched, then \texttt{buf\_len} is reset to 0;
  \item any other byte is appended to \texttt{line\_buf};
  \item if \texttt{line\_buf} would overflow 512 bytes without a newline, the
        server replies \texttt{ERR TOOLONG}, sets the \texttt{overlong} flag and
        discards bytes until the next \verb|'\n'| -- so a malicious or buggy peer
        cannot corrupt memory or desynchronise the parser.
\end{itemize}
Leftover bytes (a half command) simply stay in \texttt{line\_buf} until the next
\texttt{read()} delivers the rest.

\paragraph{Writes.} Every reply goes through \texttt{send\_all()}, which loops on
\texttt{write()} until the whole buffer is gone, retrying on \texttt{EINTR} and
advancing the offset on short writes. \texttt{SIGPIPE} is ignored so a write to a
peer that vanished returns \texttt{-1/EPIPE} instead of killing the process.

% =====================================================================
\section{Signal Handling and Shutdown}

\subsection{Server -- \texttt{Ctrl+C}}
\begin{enumerate}[nosep]
  \item \texttt{SIGINT} is caught by a handler that does exactly one
        async-signal-safe thing: \texttt{g\_stop = 1} (a
        \texttt{volatile sig\_atomic\_t}).
  \item The handler is installed \emph{without} \texttt{SA\_RESTART}, so a
        \texttt{select()} in progress returns \texttt{-1} with \texttt{errno ==
        EINTR}; the loop treats that as "re-check the flag".
  \item Seeing \texttt{g\_stop}, the loop leaves \texttt{while (!g\_stop)} and
        enters shutdown: send \texttt{SERVER SHUTDOWN} to every connected client,
        \texttt{close()} every client fd, \texttt{close()} the listening socket,
        \texttt{free()} the client array, \texttt{fclose()} the logfile.
  \item It logs \texttt{SHUTDOWN signal=SIGINT} and \texttt{CLEANUP\_DONE
        clients=0} and \texttt{return}s -- no orphan fds, no zombies (there are no
        child processes).
\end{enumerate}

\subsection{Wizard / professor -- \texttt{Ctrl+C}}
\begin{enumerate}[nosep]
  \item \texttt{SIGINT} sets \texttt{g\_stop = 1}; the client's \texttt{select()}
        loop (a 1\,s tick keeps it responsive) notices it.
  \item The client sends \texttt{APPARATE}, prints
        \texttt{[ROLE user] SENT APPARATE}, stops reading stdin, and waits for the
        reply.
  \item On \texttt{OK APPARATE} it prints \texttt{[ROLE user] DISCONNECTED
        reason=APPARATE}, \texttt{close()}s the socket and exits 0. (If the
        server is already gone it gives up after a short bounded wait.)
\end{enumerate}
The client also exits cleanly on \texttt{SERVER SHUTDOWN}
(\texttt{reason=shutdown}), on \texttt{TIMEOUT DISCONNECT}
(\texttt{reason=timeout}), and on EOF of stdin (it sends \texttt{APPARATE} first).

% =====================================================================
\section{Mandatory Test Scenarios}
\emph{Replace each placeholder with an uncut screenshot of the corresponding
terminal -- from the shell prompt through program exit. One screenshot per
terminal.}

\newcommand{\shot}[1]{%
  \begin{center}\fbox{\parbox[c][32mm][c]{0.85\linewidth}{\centering\ttfamily\small
  [ screenshot placeholder: #1 ]}}\end{center}}

\subsection{Scenario 1 -- enrollment errors, NOT\_ENROLLED, UNKNOWN, professor BREW}
Server \texttt{-n 4 -t 30}. Two clients enroll; a third tries the same username
(\texttt{ERR ENROLL name\_taken}); a fourth sends a command before enrolling
(\texttt{ERR NOT\_ENROLLED}); \texttt{FLY BROOMSTICK} $\to$ \texttt{ERR UNKNOWN
FLY}; the professor's \texttt{BREW}/\texttt{CONSUME} $\to$ \texttt{ERR
UNAUTHORIZED}; the wizard's \texttt{INSPECT}/\texttt{SCROLL} succeed; both leave
with \texttt{Ctrl+C}.
\shot{server} \shot{wizard hermione} \shot{professor dumbledore}
\shot{wizard hermione (duplicate)} \shot{third client (pre-ENROLL)}

\subsection{Scenario 2 -- BREW/CONSUME, spellbook, INSUFFICIENT, UNKNOWN ingredient}
Harry brews \texttt{MOONSTONE}/\texttt{UNICORN\_HAIR}, checks \texttt{SPELLBOOK},
over-consumes \texttt{UNICORN\_HAIR} (\texttt{ERR INSUFFICIENT\_INGREDIENTS}),
brews an unknown ingredient (\texttt{ERR UNKNOWN\_INGREDIENT}); McGonagall does
\texttt{INSPECT}/\texttt{SCROLL}/\texttt{ROSTER}.
\shot{server} \shot{wizard harry} \shot{professor mcgonagall}

\subsection{Scenario 3 -- HOGWARTS\_FULL and slot recovery}
Server \texttt{-n 3 -t 60}. Three clients fill the academy; the next two get
\texttt{ERR HOGWARTS\_FULL}; harry leaves; \texttt{neville} now fits, brews,
\texttt{ROSTER} reflects the change; everyone leaves with \texttt{Ctrl+C}.
\shot{server} \shot{harry} \shot{ron} \shot{dumbledore}
\shot{voldemort (rejected)} \shot{snape (rejected)} \shot{neville}

\subsection{Scenario 4 -- idle timeout + clean server shutdown}
Server \texttt{-n 6 -t 10}. \texttt{hermione} stays active; \texttt{harry} sends
nothing for 10\,s and is dropped with \texttt{TIMEOUT DISCONNECT}; more clients
join and query; the server gets \texttt{Ctrl+C}; \texttt{ps aux | grep hogwarts}
shows no leftover process; \texttt{cat server.log} shows the full timeline.
\shot{server + ps + cat server.log} \shot{hermione} \shot{harry (timed out)}
\shot{ron} \shot{dumbledore} \shot{mcgonagall}

% =====================================================================
\section{Challenges}

\paragraph{1. Recomputing the \texttt{select()} timeout correctly.}
A first version passed a fixed \texttt{struct timeval} to \texttt{select()}, which
meant idle clients were only noticed when \emph{some other} fd happened to be
ready. The fix was to (a) run the idle sweep \emph{before} building the fd set,
and (b) set the \texttt{select()} timeout to the minimum of
\texttt{TIMEOUT - (now - last\_active)} over all clients (and \texttt{NULL} when
there are no clients). Now the loop is guaranteed to wake up no later than the
next deadline, with no \texttt{sleep()} and no timer thread.

\paragraph{2. Message boundaries vs.\ \texttt{read()} boundaries.}
Sending two commands quickly, or a command split across packets, initially caused
either lost or merged commands. Adding a per-client \texttt{line\_buf} plus a
byte-at-a-time scan for \verb|'\n'| (with CRLF tolerance, an \texttt{ERR TOOLONG}
guard, and leftover bytes retained for the next \texttt{read()}) made the parser
robust to arbitrary fragmentation.

\paragraph{3. (Bonus) File-descriptor and slot leaks on disconnect.}
Every disconnect path -- \texttt{APPARATE}, peer hangup ($\texttt{read()} = 0$),
read error, idle timeout and global shutdown -- funnels through one
\texttt{client\_close()} that \texttt{close()}s the fd and zeroes the slot
(including the spellbook). A quick \texttt{lsof}/\texttt{ps} check after a
storm of connect/disconnect confirmed no fd growth and no zombie processes.

% =====================================================================
\section*{Appendix: Output line formats}
\addcontentsline{toc}{section}{Appendix: Output line formats}
\begin{itemize}[nosep]
  \item Server, to stdout \emph{and} the logfile, one timestamped line per
        event: \texttt{[YYYY-MM-DD HH:MM:SS] KEYWORD field=value \dots}
        (e.g.\ \texttt{SERVER\_STARTED}, \texttt{CLIENT\_CONNECTED},
        \texttt{ENROLL}, \texttt{BREW}, \texttt{CONSUME}, \texttt{INSPECT},
        \texttt{SCROLL}, \texttt{ROSTER}, \texttt{REJECTED}, \texttt{TIMEOUT},
        \texttt{CLIENT\_DISCONNECTED}, \texttt{SHUTDOWN}, \texttt{CLEANUP\_DONE}).
        Plus the one-off \texttt{Hogwarts is ready. Port: \dots} banner.
  \item Clients, to stdout: \texttt{[WIZARD <user>]} / \texttt{[PROFESSOR <user>]}
        followed by \texttt{CONNECTED}, \texttt{SENT}, \texttt{RECEIVED} or
        \texttt{DISCONNECTED}; a \texttt{> } prompt precedes each typed command.
\end{itemize}

\end{document}
